% Options for packages loaded elsewhere
\PassOptionsToPackage{unicode}{hyperref}
\PassOptionsToPackage{hyphens}{url}
\PassOptionsToPackage{dvipsnames,svgnames,x11names}{xcolor}
\documentclass[
  10pt,
]{article}
\usepackage{xcolor}
\usepackage[margin=0.8in]{geometry}
\usepackage{amsmath,amssymb}
\setcounter{secnumdepth}{-\maxdimen} % remove section numbering
\usepackage{iftex}
\ifPDFTeX
  \usepackage[T1]{fontenc}
  \usepackage[utf8]{inputenc}
  \usepackage{textcomp} % provide euro and other symbols
\else % if luatex or xetex
  \usepackage{unicode-math} % this also loads fontspec
  \defaultfontfeatures{Scale=MatchLowercase}
  \defaultfontfeatures[\rmfamily]{Ligatures=TeX,Scale=1}
\fi
\usepackage{lmodern}
\ifPDFTeX\else
  % xetex/luatex font selection
\fi
% Use upquote if available, for straight quotes in verbatim environments
\IfFileExists{upquote.sty}{\usepackage{upquote}}{}
\IfFileExists{microtype.sty}{% use microtype if available
  \usepackage[]{microtype}
  \UseMicrotypeSet[protrusion]{basicmath} % disable protrusion for tt fonts
}{}
\usepackage{setspace}
\makeatletter
\@ifundefined{KOMAClassName}{% if non-KOMA class
  \IfFileExists{parskip.sty}{%
    \usepackage{parskip}
  }{% else
    \setlength{\parindent}{0pt}
    \setlength{\parskip}{6pt plus 2pt minus 1pt}}
}{% if KOMA class
  \KOMAoptions{parskip=half}}
\makeatother
\usepackage{longtable,booktabs,array}
\setlength{\LTpre}{8pt}
\setlength{\LTpost}{8pt}
\newcounter{none} % for unnumbered tables
\usepackage{calc} % for calculating minipage widths
% Correct order of tables after \paragraph or \subparagraph
\usepackage{etoolbox}
\makeatletter
\patchcmd\longtable{\par}{\if@noskipsec\mbox{}\fi\par}{}{}
\makeatother
% Allow footnotes in longtable head/foot
\IfFileExists{footnotehyper.sty}{\usepackage{footnotehyper}}{\usepackage{footnote}}
\makesavenoteenv{longtable}
\setlength{\emergencystretch}{3em} % prevent overfull lines
\providecommand{\tightlist}{%
  \setlength{\itemsep}{0pt}\setlength{\parskip}{0pt}}
\usepackage[authoryear]{natbib}
% The approved surface labels (including W3C) pair with resolved BibTeX years.
\newcommand{\axCite}[2]{\hyperlink{cite.#1}{#2}, \citeyear{#1}}
\newcommand{\axCitet}[2]{\hyperlink{cite.#1}{#2} (\citeyear{#1})}
\usepackage{bookmark}
\IfFileExists{xurl.sty}{\usepackage{xurl}}{} % add URL line breaks if available
\urlstyle{same}
\newcommand{\doi}[1]{doi: \href{https://doi.org/#1}{\nolinkurl{#1}}}
\hypersetup{
  pdftitle={Human Authority and Institutional Intelligence in AI-Assisted Governance},
  pdfauthor={Mauricio Z. Filho},
  colorlinks=true,
  linkcolor={Maroon},
  filecolor={Maroon},
  citecolor={Blue},
  urlcolor={Blue},
  pdfcreator={LaTeX via pandoc}}

\title{Human Authority and Institutional Intelligence in AI-Assisted
Governance}
\author{Mauricio Z. Filho\\\small Axodus Founder\\\small \href{mailto:mzfshark@gmail.com}{mzfshark@gmail.com}}
\date{}

\begin{document}
\maketitle
\begin{center}\footnotesize Approved documentary identity: \texttt{AXODUS-HUMAN-AI-R1-MAJOR-REVISION-P1B-AUTHORIAL-FIX-01}\end{center}

\setstretch{1.05}
\section{Abstract}\label{abstract}

Artificial intelligence can participate in institutional processes by
assisting the formulation of proposals, analysis of alternatives,
retrieval of institutional records, contestation of claims, and
execution of bounded operational tasks. These capabilities raise a
governance question: how can institutions use distributed cognitive and
operational capabilities while preserving legitimately constituted
authority, meaningful participation, accountability, and the capacity to
challenge decisions?

This paper proposes an integrative conceptual framework for examining
AI-assisted governance in human institutions. Its normative premise is
that institutional authority must be constituted through legitimate
human institutional arrangements; technical capability or operational
performance does not independently establish such authority. The
framework distinguishes four substantive dimensions---human legitimacy,
institutional intelligence, agentic delegation, and federated
institutional autonomy---from three analytical levels: institutional
principles and authority, governance processes and human--AI
interaction, and technical implementation mechanisms.

Institutional intelligence is treated as a set of organizational
functions rather than as a claim that intelligence is generated
exclusively by artificial agents. AI assistance may be persistent,
episodic, specialized, or distributed across institutions; the capacity,
its mobilization in a process, and the resulting outcomes are distinct.
The paper examines proposal composition, deliberative analysis,
contestation, and institutional memory, together with the risks of
agenda influence, cognitive dependence, shared bias, and false
interpretive authority.

As a cross-cutting analytical device, the paper develops an
\emph{Institutional Decision Episode} that distinguishes the validity of
an institutional mandate, the conformity of a delegated action, and the
practical capacity to contest or correct an outcome. A hypothetical
procurement episode illustrates how these categories can support
different lines of inquiry. The example does not demonstrate improved
governance or operational effectiveness. The framework is a normative
and conceptual synthesis; its usefulness, effectiveness, and
applicability across institutional forms remain to be evaluated.

\textbf{Keywords:} institutional governance; human authority;
institutional intelligence; AI delegation; accountability;
contestability; socio-technical systems

\section{1. Introduction}\label{introduction}

Institutions increasingly use computational systems to support
activities that extend beyond record storage and the execution of fixed
procedures. Contemporary artificial intelligence systems can generate
text, summarize documents, identify possible relationships among
records, suggest alternatives, and operate software tools when connected
to suitable interfaces. These capabilities create opportunities to
redistribute some of the cognitive and operational work involved in
governing an organization.

The institutional implications extend beyond whether an automated
service should be allowed to perform a specific task. Governance also
depends on how problems enter an agenda, how proposals are composed,
what information reaches decision-makers, how past commitments are
interpreted, and whether affected parties can question an outcome. AI
systems may influence each of these activities even when they do not
possess a formal vote or final decision right.

This paper asks:

\begin{quote}
How can human institutions incorporate AI into their governance
processes to expand their cognitive and operational capacities while
preserving legitimate authority, meaningful participation,
accountability, and institutional autonomy?
\end{quote}

The question is normative and organizational as well as technical. A
system's ability to generate a proposal, operate a tool, or influence a
decision does not, by itself, explain whether the resulting action is
institutionally authorized. Conversely, limiting formal decision rights
to humans does not establish that human participants have sufficient
information, time, access, or practical influence to exercise those
rights meaningfully.

The framework begins from a \textbf{human-centered normative premise}:
within the institutional arrangements examined here, authority is to be
constituted through legitimate human principals and human institutions.
AI systems may contribute to institutional cognition and may perform
operational actions where authorized, but their cognitive or technical
capabilities do not independently confer institutional sovereignty. This
premise is a starting point for argument and institutional design, not
an empirical finding or a universal legal conclusion.

The paper develops an integrative account around four dimensions:

\begin{enumerate}
\def\labelenumi{\arabic{enumi}.}
\tightlist
\item
  \textbf{Human legitimacy:} the basis on which institutional authority
  is constituted, exercised, delegated, challenged, and amended.
\item
  \textbf{Institutional intelligence:} the human and AI-assisted
  cognitive functions involved in composing proposals, analyzing
  alternatives, contesting claims, and accessing institutional memory.
\item
  \textbf{Agentic delegation:} the relationship among institutional
  mandates, operational capabilities, technical permissions, actions,
  and accountability.
\item
  \textbf{Federated institutional autonomy:} how distinct institutions
  may organize their own capabilities while remaining subject to
  applicable shared commitments.
\end{enumerate}

These dimensions identify subjects of investigation. They are
complemented by three analytical levels that distinguish
\textbf{institutional authority}, \textbf{governance processes}, and
\textbf{technical mechanisms}. The two structures are not
interchangeable: one identifies substantive concerns, while the other
separates types of institutional analysis.

The paper also uses an \emph{Institutional Decision Episode} as a
cross-cutting device. It connects cognitive work before a decision to
mandate validity, action conformity, the evidence available, and the
possibility of review or correction, without treating these assessments
as interchangeable. This device integrates established concerns about
authority, delegation, accountability, and contestability; it is not
presented as the discovery of those concerns.

The framework is intended for conceptual examination across different
organizational forms. Its portability is a research proposition, not an
established result. Blockchain, smart contracts, tokenization, or any
other specific implementation technology are optional possibilities
rather than premises of the analysis.

\section{2. Theoretical Foundations and Conceptual
Method}\label{theoretical-foundations-and-conceptual-method}

\subsection{2.1 Socio-technical institutions and distributed
governance}\label{socio-technical-institutions-and-distributed-governance}

A socio-technical perspective treats institutional outcomes as arising
from interactions among people, organizational rules, practices,
information, and technical systems. \axCitet{baxter2011sociotechnical}{Baxter and Sommerville} argue
that technical design must account for organizational and social
conditions rather than treating them as external to systems engineering.
This perspective is relevant to AI-assisted governance because the
consequences of a model output depend on how it is introduced into a
workflow, who can question it, and which actors possess authority to act
on it.

Systems-oriented safety analysis likewise directs attention to
interactions and control relationships rather than isolated component
failures alone (\axCite{leveson2012engineering}{Leveson}). In the present paper, this insight is
used as a conceptual caution: failures of institutional oversight may
involve mismatches among mandates, organizational practices, available
information, and technical execution. It does not establish that the
proposed framework prevents such failures.

Polycentric governance provides a further reference for examining
institutions with multiple, interdependent centers of decision-making
(\axCite{ostrom2010beyond}{Ostrom}). Some organizations may be federated or otherwise
distributed; others may use different structures. Polycentricity is
therefore a useful comparative perspective, not a requirement that all
AI-assisted institutions share a constitutional center or federal
hierarchy.

\subsection{2.2 Legitimacy, decision rights, and delegated
action}\label{legitimacy-decision-rights-and-delegated-action}

A distinction between technically possible action and legitimate
exercise of authority is central to the proposed framework. Stone and
Mittelstadt examine the legitimacy of automated exercises of power and
caution against reducing legitimacy to other desirable properties, such
as accuracy or procedural performance (\axCite{stone2026legitimate}{Stone \& Mittelstadt}).
Their work provides an important point of comparison: compliance with a
procedure may itself require scrutiny of the authority and institutional
conditions under which that procedure operates.

Organizational research also examines how decision rights are
distributed among human and AI participants. \axCitet{sebastian2026decision}{Sebastian et al.}
discuss decision rights for AI in an organizational context. \axCitet{engin2026human}{Engin} examines configurations of decision authority, process autonomy,
and accountability. These approaches overlap materially with the
distinctions pursued here. The present paper does not claim to have
originated the separation between authority and operational autonomy.

The more specific question developed below concerns how an institutional
episode can be examined when its mandate, execution, and review
conditions do not align. This is proposed as an organizing synthesis.
Whether it produces a substantive contribution beyond existing
frameworks remains a matter for further comparison.

\subsection{2.3 Human control, reviewability, and
contestability}\label{human-control-reviewability-and-contestability}

Meaningful human control cannot be assumed merely because a person
appears somewhere in a decision process. \axCitet{santonidesio2018meaningful}{Santoni de Sio and van den Hoven} analyze conditions under which autonomous systems can
remain connected to human reasons and responsibility. Their account is
relevant to the distinction between nominal human involvement and a
substantively meaningful relationship between human agency and system
behavior.

Research on the oversight of automated decisions also questions policies
that treat a human review step as sufficient by default (\axCite{green2022flaws}{Green}).
The relevance here is not that every human reviewer will fail, but that
oversight depends on institutional conditions: access to information,
practical capacity, authority to intervene, and an appropriate
understanding of the task.

Accountability and contestability scholarship adds a complementary
concern. \axCitet{kroll2017accountable}{Kroll et al.} argue that transparency alone is not an
adequate account of algorithmic accountability. \axCitet{cobbe2021reviewable}{Cobbe et al.}
examine the reviewability of automated decision-making as a
socio-technical process. \axCitet{alfrink2023contestable}{Alfrink et al.} develop contestability as
a design concern. Together, these works support examining how decisions
can be reconstructed, questioned, and reconsidered; they do not
establish that a particular record or review procedure necessarily
produces an effective remedy.

\subsection{2.4 Related work and the status of the
contribution}\label{related-work-and-the-status-of-the-contribution}

This article offers a bounded integrative synthesis. HAIG examines the
distribution of decision authority, process autonomy, and accountability
in human--AI arrangements (\axCite{engin2026human}{Engin}). Reviewable automated
decision-making situates records and accountability across the
sociotechnical process (\axCite{cobbe2021reviewable}{Cobbe et al.}), while contestable AI by
design specifies features and practices supporting intervention and
challenge (\axCite{alfrink2023contestable}{Alfrink et al.}). \axCitet{sebastian2026decision}{Sebastian et al.} address
organizational decision rights in relation to ambiguity and risk.

The Institutional Decision Episode (IDE) makes selected distinctions
explicit in a common record for one decision or action, or for linked
episodes. Its potential incremental utility is clearest where an analyst
must distinguish a change of institutional competence from a change in
operational permissions, or deliberative approval from subsequent
execution. The B4 reconstruction illustrates how the scheme can preserve
these distinctions. It does not show that other approaches cannot
represent them or that the scheme improves analysis in practice. Where
an existing framework answers the question at the required resolution,
the IDE may add no value.

\subsection{2.5 Conceptual method}\label{conceptual-method}

This article combines an integrative conceptual synthesis, a structured
comparison with selected scholarly accounts, and analytical
reconstruction of hypothetical episodes. It asks how AI-assisted
cognitive and operational work can be distinguished from institutional
authority, and under what conditions an episode representation might
clarify that distinction. Sources were selected for their relevance to
organizational cognition, authority, delegation, reviewability, or
contestability. Targeted bibliographic checks used title, author,
identifier, and topic searches across public scholarly sources. No
complete historical search or screening log is available. Coverage is
purposive, and access ranges from reviewed passages in full texts to
abstracts and metadata.

Comparisons distinguish unit of analysis, analytical purpose, explicitly
represented relationships, evidence requirements, and cases of
redundancy. The procurement illustration separates three diagnostic
concerns. Scenario A illustrates the boundary between bounded
authorization and external execution; Scenario B illustrates linked
local and federative decisions. The reconstruction of B4 uses the same
stipulated record and reports missing evidence rather than supplying it.
These materials demonstrate initial applicability only. No empirical
validation or independent inter-analyst replication was performed.
Source selection, uneven access, and the use of a constructed example
limit inference. Future assessments require fixed records, explicit
comparators, independent coding, and transparent treatment of null
results and disagreements.

This is not a systematic review, meta-analysis, controlled experiment,
or formal proof. Proposed evaluation designs in Section 9 require
appropriate access, institutional approval, and ethical review before
any future execution.

\section{3. Legitimate Authority and Institutional Decision
Rights}\label{legitimate-authority-and-institutional-decision-rights}

\subsection{3.1 Human-centered authority as a normative
premise}\label{human-centered-authority-as-a-normative-premise}

``Intelligence may be distributed; authority must be constituted'' is a
premise of the institutional model examined here, not a universal theory
of moral or political legitimacy. An actor's technical capability or
epistemic contribution does not by itself confer a right to decide.

Formal constitution and recognition identify who is treated as holding
authority. Internal institutional validity asks whether that actor had
competence and followed applicable rules. Legal validity, where
relevant, is a separate question under the applicable jurisdiction.
Social acceptance concerns actual beliefs or acquiescence; normative
legitimacy concerns whether the exercise of power is justified to those
affected. These assessments can diverge. A formally authorized decision
can remain normatively contestable, including when affected parties lack
representation, meaningful participation, rights protection, or a
practical opportunity to challenge it. \axCitet{stone2026legitimate}{Stone and Mittelstadt}
distinguish descriptive acceptance from the normative question of who
may permissibly exercise power.

A recognized institutional principal or body is an actor or collective
body treated, under the arrangements being examined, as holding or
exercising a specified decision right. Recognition does not settle
whether the underlying rule or distribution of power is legitimate. The
IDE records the asserted basis of authority and the evidence and avenues
for its challenge; it supplies no universal test of legitimacy. The
framework does not decide universal questions of AI legal personhood,
rights, or standing, which remain open and context dependent.

\subsection{3.2 Capability, permission, factual agency, and
authority}\label{capability-permission-factual-agency-and-authority}

Four concepts should be distinguished.

{\def\LTcaptype{none} % do not increment counter
\setlength{\parskip}{0pt}\begin{longtable}[]{@{}
  >{\raggedright\arraybackslash}p{(\linewidth - 4\tabcolsep) * \real{0.3333}}
  >{\raggedright\arraybackslash}p{(\linewidth - 4\tabcolsep) * \real{0.3333}}
  >{\raggedright\arraybackslash}p{(\linewidth - 4\tabcolsep) * \real{0.3333}}@{}}
\toprule\noalign{}
Concept & Working meaning & What it does not establish by itself \\
\midrule\noalign{}
\endhead
\bottomrule\noalign{}
\endlastfoot
\textbf{Technical capability} & What a system can materially compute,
generate, access, or execute in an environment. & That any action is
permitted or legitimate. \\
\textbf{Configured enablement} & An identifiable technical or procedural
mechanism that enables an actor to invoke an operation or register an
act. & That the mechanism exists merely because a rule, mandate, or
completed vote exists; or that it reflects a valid institutional
mandate. \\
\textbf{Operational agency} & The factual capacity to act or influence
events, including by initiating tool calls or shaping available
information. & That the action was authorized or that the actor holds
legitimate institutional authority. \\
\textbf{Institutional authority} & A decision right or competence
constituted or recognized under the institutional arrangements being
examined. & That the right is technically enforceable or uncontested in
every relevant context. \\
\end{longtable}
}

In the IDE, enablement is recorded as technical (T\_tech) or procedural
(T\_proc) only when a distinct mechanism is evidenced. Where relevant
but not shown in examined records it is NOT\_EVIDENCED; where
applicability is unclear it is UNKNOWN. NOT\_APPLICABLE concerns a
specified type or stage of act, not proof that a mechanism does not
exist elsewhere.

Operational agency is a \textbf{factual} concept in this model. A system
may act outside its mandate, use an overly broad technical permission,
or perform an action without any valid institutional authorization. It
may also influence a decision through agenda selection or presentation
of information without possessing an explicit execution permission.

The distinction matters because an access-control mechanism may permit
an action that the institution did not validly authorize. It also
matters in the opposite direction: an institution may intend to
authorize an action, while the technical system is unable or incorrectly
configured to perform it. In neither case does the observed technical
state settle all institutional questions.

An action lacking a valid mandate should therefore be classified
initially as an \textbf{authority or authorization concern} under the
framework. Whether it is institutionally invalid, legally void,
voidable, capable of ratification, or binding on an external party
depends on the applicable rules and circumstances. The model does not
presume a universal legal consequence.

\subsection{3.3 Candidate questions about delegation and retained
authority}\label{candidate-questions-about-delegation-and-retained-authority}

The human-centered premise does not require direct human approval of
every operational act. An institution may wish to authorize bounded
automation while retaining authority over the mandate, its limits, and
the procedures for review or amendment.

Which powers may be delegated, and which must remain reserved to human
institutional bodies, is not resolved here. Five categories deserve
particular institutional examination:

\begin{enumerate}
\def\labelenumi{\arabic{enumi}.}
\tightlist
\item
  \textbf{Constitutive and amendment powers:} authority to establish or
  revise foundational rules.
\item
  \textbf{Allocation of institutional competence:} authority to create
  roles, grant mandates, or alter their scope.
\item
  \textbf{Participation and standing:} authority to determine
  membership, voting rights, or access to institutional procedures.
\item
  \textbf{Adjudication and remedy:} authority to issue binding
  dispositions on disputes or challenges.
\item
  \textbf{Emergency powers and exceptions:} authority to suspend
  ordinary procedures or exercise exceptional powers.
\end{enumerate}

These are \textbf{candidate areas for determination}, not a universally
established list of nondelegable powers. Their treatment may differ by
institution, applicable law, risk, reversibility, and the nature of the
decision. The distinction between authority over a standing mandate and
approval of individual actions within that mandate remains central to
the open question of Human Finality.

\subsection{3.4 Responsibility without automatic
individualization}\label{responsibility-without-automatic-individualization}

Accountability requires more than assigning a name to an outcome.
Institutional responsibility may involve a collective body that
established a mandate, people who configured a process, operators who
supervised it, and reviewers or other parties with relevant duties.
Which roles bear responsibility depends on the institution's rules and
the facts of the episode.

The framework therefore seeks \textbf{attributable institutional
responsibility} rather than presuming that every outcome must be
assigned to one individual. Attribution for governance analysis is also
distinct from a determination of legal liability.

\section{4. AI-Augmented Institutional
Intelligence}\label{ai-augmented-institutional-intelligence}

\subsection{4.1 Institutional intelligence as a governance
function}\label{institutional-intelligence-as-a-governance-function}

Institutional intelligence denotes a governance arrangement's capacity
to formulate issues, examine alternatives and evidence, receive and
respond to challenges, and retrieve and correct records relevant to
institutional judgment. The unit of analysis is a bounded governance
arrangement, identified by its decision domain, participants, routines,
and period of observation. Individuals, groups, records, and technical
artifacts can sustain that capacity; their outputs are not measures of
it by themselves. Work on bounded rationality, organizational
information processing, and distributed cognition motivates examining
these arrangements rather than attributing the capacity to a single
actor (\axCite{simon1978rational}{Simon}; \axCite{galbraith1974organization}{Galbraith}; \axCite{hollan2000distributed}{Hollan et al.}).

Composition, analysis, contestation, and memory are four overlapping
governance functions, not an exhaustive taxonomy of cognition. Attention
shapes what each function makes visible; coordination connects
participants and tasks; deliberation brings reasons and objections into
collective consideration; learning concerns changes in routines over
time. Retrieving a record does not establish that the institution has
learned from it (\axCite{walsh1991memory}{Walsh \& Ungson}; \axCite{levitt1988learning}{Levitt \& March}).

AI assistance can redistribute cognitive work, but its effects depend on
the task, sources, opportunities for correction, and division of
responsibilities. Combining more outputs does not establish
complementary judgment or independent scrutiny (\axCite{gonzalez2026teaming}{Gonzalez et al.}; \axCite{vaccaro2024combinations}{Vaccaro et al.}). A useful contribution to judgment does not confer
authority over the resulting decision. Assistance may be continuous,
occasional, specialized, centralized, or distributed. Benefits are
possibilities and hypotheses, not established outcomes.

\subsection{4.2 Proposal composition and institutional
participation}\label{proposal-composition-and-institutional-participation}

Proposal composition is an institutional process because it determines
how a concern becomes an object of collective attention and potential
decision. A proposal must normally be intelligible to the relevant
participants and connected to a process through which they can consider
it.

AI assistance may help a participant clarify an objective, identify
missing information, organize arguments, or locate potentially relevant
rules. It may also offer alternative formulations or explain what
further questions a proposal leaves unresolved. These functions could
lower some drafting barriers for some participants; whether they do so,
and for whom, requires evaluation.

Composition is not a neutral transcription of a participant's intention.
A system may suggest objectives, select a problem definition, omit
alternatives, or embed assumptions about what is feasible. The
participant and the institution therefore need ways to distinguish the
participant's stated purpose from AI-generated additions or reframing.

Wider access to composition assistance may support participation, but
access alone does not establish equal understanding, institutional
standing, influence over an agenda, or political power. Differences in
resources, experience, accessibility, and control over the tools may
persist or change form.

\subsection{4.3 Analysis and deliberative
assistance}\label{analysis-and-deliberative-assistance}

Institutional deliberation often requires participants to compare
alternatives under incomplete information. AI systems may help summarize
source materials, identify apparent inconsistencies, assemble questions
about consequences, or present competing interpretations for
examination.

The institutional role of analysis is to support deliberation, not to
convert a model-generated assessment into a binding disposition. A
useful analysis should allow participants to distinguish:

\begin{enumerate}
\def\labelenumi{\arabic{enumi}.}
\tightlist
\item
  the source records or evidence consulted;
\item
  the interpretation or inference generated from those records;
\item
  assumptions and unresolved uncertainties;
\item
  alternatives that were considered or omitted; and
\item
  the body or role competent to decide the relevant issue.
\end{enumerate}

These are proposed conditions for reviewable assistance. They are not a
claim that every system can provide a complete explanation or that
source links alone guarantee correctness.

Analytical assistance also has distributional consequences. The same
summary can help participants locate relevant material while
concentrating influence in whoever selects the sources, configures the
task, or controls the presentation of alternatives. Institutions must
therefore be able to examine not only the answer produced, but also the
process by which the question and available options were framed.

\subsection{4.4 Adversarial analysis and institutional
contestation}\label{adversarial-analysis-and-institutional-contestation}

A proposal can be clearly composed and still rest on a mistaken
assumption, incomplete evidence, conflicting obligation, or contested
interpretation. An adversarial function can assist by seeking
counterarguments, overlooked stakeholders, possible conflicts, and
alternative readings of the evidence.

This function should be distinguished from \textbf{institutional
adjudication}. An AI-generated challenge may identify a possible
conflict, but the consequence of that challenge depends on the
institution's rules and the authority of the relevant human body. A
technical mechanism might implement an applicable prohibition where the
institution has validly authorized that mechanism. An AI interpretation
of an ambiguous rule should not automatically be treated as equivalent
to an authoritative adjudication.

Neither a separate agent label nor an apparent disagreement between two
systems establishes independence. Systems may share models, providers,
source collections, incentives, or assumptions. Assessing adversarial
scrutiny therefore requires examining the relevant sources of dependence
and the practical opportunity for human participants or independent
reviewers to challenge the analysis.

Contestatory assistance is also distinct from a formal right of appeal.
An institution may provide criticism before a decision while offering no
effective route to challenge the resulting action. Conversely, a review
procedure may exist without the resources needed to use it. The
framework treats these as related but separate institutional conditions.

\subsection{4.5 Active institutional
memory}\label{active-institutional-memory}

Institutions retain decisions, policies, mandates, commitments,
exceptions, and other records over time. The difficulty is not only
storing those records, but making relevant records accessible when an
institution considers a new proposal or reviews an earlier action.

AI-assisted memory may help participants locate a prior decision,
identify a possible relationship between records, or retrieve materials
relevant to a current question. Provenance models such as PROV-DM
provide a vocabulary for describing entities, activities, and relations
among records (\axCite{w3c2013provdm}{World Wide Web Consortium {[}W3C{]}}). Their use
does not by itself establish that a record is authentic, complete,
current, or authoritative.

A distinction is required between:

\begin{itemize}
\tightlist
\item
  an institutional record;
\item
  information about the record's provenance, status, and possible
  supersession; and
\item
  an AI-generated interpretation of what that record means for a current
  decision.
\end{itemize}

A system can retrieve an obsolete decision, omit a later amendment,
disclose information that should not be shared, or present a
model-generated summary as though it were a binding precedent. Effective
use of institutional memory therefore depends on governance of sources,
access, status, interpretation, and correction. The appropriate
mechanisms will vary by institution and by the sensitivity of the
records involved.

\subsection{4.6 Cognitive assistance and legitimate institutional
decision-making}\label{cognitive-assistance-and-legitimate-institutional-decision-making}

The four functions may shape a decision before any formal vote or
authorization occurs. Composition can frame the problem; analysis can
select evidence and alternatives; contestation can determine which
objections are visible; and memory can influence which prior commitments
appear relevant.

This creates a distinction between \textbf{formal decision authority}
and \textbf{practical influence over a decision}. A human body may
retain the formal right to decide while relying heavily on a system that
shapes the available information. That possibility does not establish
that such a transfer of influence has occurred in a particular
institution. It does, however, define an institutional question that
cannot be answered solely by asking whether the system possesses a vote.

Meaningful human participation may require more than the presence of an
approval step. Candidate conditions include access to relevant evidence,
time and competence appropriate to the task, the ability to request
alternatives or challenge interpretations, and practical authority to
reject, suspend, or escalate an action. The adequacy and proportionality
of these conditions remain subject to institutional and empirical
examination.

\subsection{4.7 Risks and competing
interpretations}\label{risks-and-competing-interpretations}

Four risks are particularly relevant to the proposed cognitive
functions.

\textbf{Agenda capture:} assistance may privilege particular problems,
sources, or alternatives. Influence can arise through selection and
framing even without formal voting power.

\textbf{Cognitive dependence:} participants may rely on generated
explanations without independently checking matters that require their
judgment. The conditions under which assistance improves or weakens
independent judgment remain an empirical question.

\textbf{Epistemic monoculture:} systems presented as independent may
rely on similar models, sources, or assumptions. This can produce
correlated omissions or reinforce one interpretation across an
institution.

\textbf{False interpretive authority:} generated summaries or apparent
rule interpretations may be treated as binding decisions or verified
precedents, even when they are provisional or wrong.

Possible responses include accessible source records, opportunities for
alternative analysis, disclosure of relevant dependencies, and review
procedures proportionate to consequences. These are design options for
examination. None guarantees independent judgment, neutrality, or
effective contestation.

\section{5. The Three-Level Governance
Framework}\label{the-three-level-governance-framework}

The three levels below distinguish different analytical questions. They
do not prescribe organizational departments, software services, or a
single institutional hierarchy.

\subsection{5.1 Level I --- Institutional principles and legitimate
authority}\label{level-i-institutional-principles-and-legitimate-authority}

Level I concerns the basis on which an institution recognizes decision
rights and allocates competencies. Its questions include:

\begin{itemize}
\tightlist
\item
  Who may establish or amend a mandate?
\item
  Which body has competence over the matter?
\item
  What procedures and constraints apply?
\item
  How can authority itself be challenged or revised?
\end{itemize}

A Level I analysis identifies an asserted authority basis. It does not
assume that the institution, its founding rules, or every decision made
under them is beyond challenge.

\subsection{5.2 Level II --- Governance processes and human--AI
interaction}\label{level-ii-governance-processes-and-humanai-interaction}

Level II concerns the processes through which institutional matters are
formulated, analyzed, deliberated, authorized, delegated, monitored, and
contested.

The cognitive functions discussed in Section 4 principally operate at
this level, though they depend on institutional authority and technical
arrangements. A process can include human participants, AI assistance,
standing operational mandates, review bodies, or other arrangements
suited to the institution.

A formal process may satisfy its written requirements while remaining
inaccessible, poorly informed, or practically unchallengeable. Thus,
procedural compliance and meaningful participation should be examined
separately.

\subsection{5.3 Level III --- Technical infrastructure and
implementation
mechanisms}\label{level-iii-technical-infrastructure-and-implementation-mechanisms}

Level III concerns the mechanisms through which institutional processes
may be represented or supported. Possible mechanisms include identity
systems, permissions, workflow tools, evidence stores, audit records,
agent runtimes, cryptographic credentials, smart contracts, or
distributed ledgers.

No listed mechanism is a universal prerequisite. A technical permission
can help implement an institutional rule but may be broader, narrower,
outdated, or otherwise inconsistent with the mandate it represents. A
technical record can support reconstruction but does not by itself
establish the truth of its contents or the legitimacy of the action
recorded.

\subsection{5.4 Relationships across the
levels}\label{relationships-across-the-levels}

The levels can be used to examine whether an institutional claim at one
level is adequately connected to another.

{\def\LTcaptype{none} % do not increment counter
\setlength{\parskip}{0pt}\begin{longtable}[]{@{}
  >{\raggedright\arraybackslash}p{(\linewidth - 2\tabcolsep) * \real{0.5000}}
  >{\raggedright\arraybackslash}p{(\linewidth - 2\tabcolsep) * \real{0.5000}}@{}}
\toprule\noalign{}
Relationship & Question for analysis \\
\midrule\noalign{}
\endhead
\bottomrule\noalign{}
\endlastfoot
\textbf{Authority → process} & Is the process grounded in a decision
right or mandate recognized under the relevant institutional
arrangements? \\
\textbf{Process → technical mechanism} & Does the configured mechanism
represent the scope and conditions of the process it is intended to
support? \\
\textbf{Technical action → institutional record} & What can available
evidence establish about what happened, and what remains uncertain? \\
\textbf{Record → review} & Can the appropriate participants or reviewers
obtain and interpret relevant evidence in time to exercise their
roles? \\
\textbf{Review → correction} & What forms of suspension, revocation,
correction, reversal, or remedy are available, and to whom? \\
\end{longtable}
}

These relationships are \textbf{proposed analytical requirements}, not
verified properties of any implementation. They should be applied
proportionately: the evidence and review needs of a routine,
low-consequence task need not be identical to those of a consequential
institutional decision.

The framework proposes that institutions examine authority provenance,
delegation boundaries, attributable responsibility, reviewability,
proportionality, and the availability of correction. Their exact form,
exceptions, and enforcement remain questions for the relevant
institution and for future evaluation.

\section{6. Accountable Delegation and Institutional Decision
Episodes}\label{accountable-delegation-and-institutional-decision-episodes}

\subsection{6.1 The institutional decision
episode}\label{the-institutional-decision-episode}

An Institutional Decision Episode is a bounded analytical representation
of a decision or delegated action. The canonical representation remains:

\[
\mathcal{E}=\langle P,R,M,C,T,A,E,Q,D\rangle
\]

Here P is an asserted principal or body; R the applicable rules and
procedures; M the mandate or claimed authorization; C relevant factual
capabilities; T evidenced configured technical or procedural enablement;
A the focal action, proposed action, or observed event; E the available
evidence; Q conditions for challenge or review; and D the actual or
possible institutional disposition.

A focal act enters the relevant procedure and ends with its immediate
disposition or the observation cutoff. Where principal, competence,
mandate, or action stage changes, separately attributable episodes may
be linked rather than merged. Joint principals and contested authority
should be recorded as such. Distinguish proposal, approval, operational
authorization, submission, external acceptance, and execution when
relevant; do not infer a later event from an earlier one.

Mandate Validity Assessment (MVA) asks whether the asserted mandate was
constituted or attributed by a competent body under applicable rules.
Report VALID, DEFECTIVE, CONTESTED, or UNDETERMINED. Independently,
mandate conformity concerns whether action A fits the scope and
conditions of mandate M under applicable rules R. We denote this
substantive relation as MC(A,M,R). Any reported conformity
classification is an evidence-supported assessment: the analyst must
identify the records E, applicable assumptions, and unresolved
interpretations on which the classification depends. For MC report
CONFORMING, NONCONFORMING, CONTESTED, or UNDETERMINED. Thus, the rules
determine the relevant normative standard, while the evidence supports
or limits the analyst's conclusion. An action can fit the nominal text
of an invalid or contested mandate; a valid mandate does not establish
action conformity. Neither a conforming classification nor a valid
mandate establishes legal effect or normative legitimacy.

For each element, an analyst identifies the following and records
unknown or contested interpretations: P, the asserted principal,
including any joint or competing principals; R, constitutive,
procedural, substantive, and where applicable legal rules; M, its
issuing body, scope, conditions, and authority; C, factual cognitive and
operational capabilities; T, a separately evidenced enablement mechanism
and its technical or procedural type; A, one focal act and its stage; E,
the sources supporting or contradicting the entry; Q, standing, access,
time, cost, independence, and review capacity; and D, the immediate
disposition and any pending later action. A record of the episode should
identify missing sources rather than completing fields from assumptions.

For each material assertion identify the record, scope, and whether it
was examined. Distinguish NOT\_EVIDENCED, EVIDENCED\_ABSENT,
NOT\_APPLICABLE, UNKNOWN, and CONTESTED. An absence claim requires
evidence appropriate to that claim; a prohibition does not prove an
event did not occur. T\_proc requires an identifiable procedural
enablement distinct from R and M; a recorded vote alone does not
establish it. Distinguish formal avenues of challenge from the practical
ability to use them, and actual dispositions from pending or prospective
ones.

The representation is a conceptual aid, not a required software
protocol, a complete account of institutional legitimacy, or evidence
that any Axodus system implements the model. Initial application to
stipulated material does not validate inter-analyst reproducibility.

\subsection{6.2 Three diagnostic
categories}\label{three-diagnostic-categories}

The framework distinguishes three candidate categories. They may occur
separately or together.

\subsubsection{6.2.1 Authority validity
concern}\label{authority-validity-concern}

An action may be technically possible, permitted by a configured system,
and consistent with the written terms of a purported mandate, while the
mandate's institutional basis is defective or disputed.

The analytical question is whether the body issuing the mandate had the
relevant competence and followed applicable constitutive procedures. A
technical permission or execution record does not answer that question
by itself.

The consequences of a defect depend on the applicable institutional and
legal rules. Possible responses may include suspension, further review,
correction of the authority record, or consideration of ratification
where permitted.

\subsubsection{6.2.2 Delegated-action conformity
concern}\label{delegated-action-conformity-concern}

A mandate may be validly constituted while a particular action exceeds
its scope or conditions. This can occur, for example, when a technical
permission is broader than the mandate or when an action violates a
condition that the technical system does not represent.

The analytical question is whether the action conforms to the mandate,
not whether the mandate itself was validly created. Depending on the
facts, review may concern configuration, operational supervision,
interpretation of scope, or other institutional responsibilities.

\subsubsection{6.2.3 Contestability capacity
concern}\label{contestability-capacity-concern}

A mandate may be valid and an action may conform to it, while an
affected participant lacks timely evidence, access to an appropriate
reviewer, or a practical opportunity to challenge or seek correction.

The analytical question is whether relevant review or corrective
capacity exists in practice. A contestability concern does not
automatically make the original action invalid. Nor does it establish
that reversal is feasible. The institution may need to distinguish
review, prospective correction, reversal where possible, and other
remedies.

\subsection{6.3 HDE-01: A hypothetical procurement
episode}\label{hde-01-a-hypothetical-procurement-episode}

The following example is \textbf{fictional and nonempirical}. It
illustrates how the three categories can be applied to a common episode;
it does not establish effectiveness, legal consequences, or operational
performance.

Assume a fictional organization whose rules permit a properly
constituted committee to delegate routine purchasing within a defined
category and spending ceiling. An AI-supported procurement service
possesses technical credentials that can be configured to submit
purchase orders. The institution also provides a review process for
affected participants.

\subsubsection{Variant A --- Authority
validity}\label{variant-a-authority-validity}

Assume that the rules require a quorum to issue a purchasing mandate,
but the mandate was recorded without that quorum. The procurement
service submits an order that falls within the amount configured in the
technical system and within the nominal terms of the recorded mandate.

The episode raises an \textbf{authority validity concern}: technical
permission and nominal conformity do not determine whether the mandate
was validly constituted. A review might consider the mandate's status,
the roles involved in recording and configuring it, and whether
suspension or another disposition is appropriate. Ratification or
cancellation would be possible responses \textbf{only if} the applicable
rules and circumstances permit them.

\subsubsection{Variant B --- Delegated-action
conformity}\label{variant-b-delegated-action-conformity}

Assume that a valid mandate permits purchases up to \textbf{\$25,000} in
an approved category. A technical configuration error permits a
\textbf{\$29,000} order, and the service submits it.

The episode raises a \textbf{delegated-action conformity concern}: the
mandate is assumed valid, but the specific order exceeds its stated
ceiling. Relevant inquiry could include how the technical permission
came to exceed the mandate, what supervision or checks applied, and what
actions remain available. Responsibility cannot be assigned solely from
the fact that the configured permission was too broad.

\subsubsection{Variant C --- Contestability
capacity}\label{variant-c-contestability-capacity}

Assume that the mandate is valid, the technical permission matches it,
and the service submits a \textbf{\$22,000} order in the approved
category. An affected participant identifies a potentially material
concern, but cannot obtain the relevant evidence in time or access an
effective review channel.

The episode raises a \textbf{contestability capacity concern}. The
action's conformity to the mandate does not settle whether the affected
participant had a meaningful opportunity to question it. A review might
examine standing, evidence access, reviewer independence, timing, and
available dispositions. The purchase is not automatically invalid merely
because a concern was raised, and reversal may or may not be possible.

\subsection{6.4 High-cadence operational delegation and external
irreversibility}\label{high-cadence-operational-delegation-and-external-irreversibility}

Whereas HDE-01 isolates the three diagnostic concerns in a discrete
purchasing episode, some delegated operations occur at a cadence where
human approval of each individual action is impractical. Retaining human
institutional authority then depends on constituting a standing mandate
over a class of bounded actions, rather than manually authorizing every
operational event.

Consider a fictional, human-governed Organization Alpha with its own
treasury. Assume that Alpha's governing board authorizes a competent
investment committee to establish a time-limited mandate for an
automated trading strategy within a board-approved capital allocation.
Before adopting the mandate, committee members may use AI assistance to
compare alternatives, retrieve prior decisions, and stress-test
assumptions, while retaining authority to reject the program. The
mandate specifies permitted instruments and venues, capital allocation,
gross leverage, single-instrument net exposure ceilings, monitoring
duties, and rules for suspension, amendment, revocation, and review.

Once constituted, operational execution follows a bounded sequence:

\textbf{Strategy / quantitative analysis → proposal → portfolio
evaluation → risk decision → pre-execution authorization → trading
intent → execution intent → execution harness → exchange boundary.}

An automated strategy can generate a candidate proposal, but proposal
generation confers no permission to act. Portfolio evaluation checks
capital allocation and aggregate exposure, and a separate risk decision
evaluates the candidate against mandate limits. Pre-execution
authorization then verifies that the proposal, portfolio and risk
decisions, mandate revision, and current operating state support issuing
a single-use trading intent, which is translated into an execution
intent for the harness to submit across the external exchange boundary.
Verifying mandate revisions and operating state at this boundary is
treated here as a design requirement of the hypothetical arrangement,
not as a verified universal property of all production systems.

Suppose the mandate caps Alpha's net long notional exposure in one
instrument at \textbf{500,000 currency units} on a declared common
valuation snapshot, distinct from gross portfolio exposure, total
capital allocation, and leverage. If current net long exposure is
\textbf{420,000 units} and a candidate proposal would add
\textbf{150,000 units} with no offsetting orders, the resulting
\textbf{570,000-unit} net exposure exceeds the ceiling by \textbf{70,000
units}. Risk rejects the proposal and pre-execution authorization issues
no trading intent. That technical rejection is an operational
disposition, not an institutional adjudication of mandate validity or
remedies.

Similarly, revoking the mandate terminates prospective authority, and
applicable admission and enforcement boundaries must reject subsequent
unauthorized actions. For inflight operations, the harness may transmit
cancellation requests for outstanding external orders, but cancellation
may fail or race with execution. Orders already matched at the external
venue cannot be undone by internal revocation, leaving residual exposure
that must be reconciled and managed under a separate valid
authorization. This fictional illustration does not establish universal
runtime enforcement, a verified cross-system authority chain, guaranteed
order cancellation, or improved economic performance.

\subsection{6.5 What the examples can and cannot
establish}\label{what-the-examples-can-and-cannot-establish}

Together, HDE-01 and the high-cadence operational illustration isolate
different analytical questions within institutional settings. They
suggest that review of mandate formation, review of action scope, and
review of contestability may call for different evidence and involve
different institutional roles.

The examples do not demonstrate that the model is more effective than
existing frameworks, that institutions will implement the distinctions,
or that any particular remedy is legally available. Responsibility and
corrective responses remain conditional on facts, institutional rules,
and applicable legal context.

The categories also are not mutually exclusive in practice. A defective
mandate, an out-of-scope action, and inaccessible review may occur in
the same episode. Their separation is useful only to the extent that it
helps an analyst avoid treating evidence of one condition as proof of
the others.

\section{7. Institutional Variation and Federated
Autonomy}\label{institutional-variation-and-federated-autonomy}

\subsection{7.1 Conceptual applicability across institutional
forms}\label{conceptual-applicability-across-institutional-forms}

The proposed framework is intended to support analysis across different
institutional forms. Its starting point is not that all organizations
possess the same authority structure, participation rules, or technical
infrastructure. Rather, it asks whether the distinctions among
institutional authority, governance processes, and technical mechanisms
can help identify relevant questions in each setting.

This is a proposition about \textbf{conceptual applicability}. It is not
evidence that the framework has been validated across institutions or
that the same arrangements should be adopted in every context.
Organizational constitution, applicable law, distribution of decision
rights, resources, and opportunities for challenge may differ
substantially.

The examples below are hypothetical configurations. They illustrate how
the framework might be used; they do not define necessary
characteristics of the institutions described.

\subsubsection{Conventional organizations and
enterprises}\label{conventional-organizations-and-enterprises}

Consider an organization in which a board establishes a purchasing
policy, a management function operates within a delegated budget, and a
technical service assists with supplier selection. The institution may
use AI to help compose purchasing proposals, retrieve relevant
contracts, identify apparent conflicts, or execute certain routine
actions under a standing mandate.

The three analytical levels help separate the authority to establish the
purchasing mandate, the processes through which purchasing decisions are
prepared and reviewed, and the technical permissions through which
orders are transmitted. The arrangement may differ in another
enterprise: some decisions may be reserved to a committee, others to
individual officers, and others governed by standing operational rules.

The use of AI does not determine which allocation is legitimate. Nor
does a technically successful order, by itself, establish that the
applicable institutional mandate was valid or followed. A further
question is whether affected participants can obtain the information
needed to raise a concern and reach an appropriate reviewer.

\subsubsection{Cooperatives and participatory
institutions}\label{cooperatives-and-participatory-institutions}

A cooperative or participatory institution may organize proposal rights
and decision procedures among members, representatives, or other
recognized bodies. Its specific voting and participation arrangements
depend on its constitution and applicable rules; no single voting
formula is assumed here.

AI-assisted composition may help participants formulate proposals in a
common format or locate relevant institutional records. Analysis and
contestation functions may also help participants compare alternatives
and identify uncertainties. These possibilities do not establish that
participation becomes more equal: members may differ in access to tools,
time, institutional standing, language, expertise, or influence over the
agenda.

For this setting, the framework directs attention to whether assistance
expands the range of proposals participants can meaningfully consider,
or whether common tools concentrate interpretive power in those who
configure them.

\subsubsection{Federated and multi-organization
structures}\label{federated-and-multi-organization-structures}

A federation or other multi-organization arrangement may distribute
authority among a common institution and constituent organizations. The
relevant allocation can arise from a constitution, agreement, statute,
or another institutional instrument. The degree of local autonomy, the
basis of membership, and the authority of common bodies vary by
arrangement; voluntary participation and uniform local powers are not
assumed.

AI-assisted institutional intelligence may operate at more than one
level. A local organization might use an assistant to relate a proposed
policy to its own records, while a common body might use another process
to examine obligations shared across the federation. A possible conflict
between local discretion and a common commitment would then require
interpretation under the applicable rules and by the body competent to
resolve it.

Ostrom's account of polycentric governance offers a conceptual reference
for examining multiple, interdependent centers of decision-making. It
does not establish that every federated arrangement is polycentric, or
that an AI-assisted federation will preserve effective local autonomy
(\axCite{ostrom2010beyond}{Ostrom}).

\subsubsection{DAOs and tokenized governance
arrangements}\label{daos-and-tokenized-governance-arrangements}

Some organizations may use decentralized or tokenized mechanisms for
aspects of proposal submission, voting, or execution. These arrangements
differ in how they recognize participation, allocate decision rights,
and connect collective decisions to technical systems. A particular
token, voting formula, smart contract, or account structure is not
assumed to be necessary for every DAO.

The framework can be used to ask whether an account or technical
mechanism acts under an institutionally recognized mandate, how the
relevant decision procedure is justified, and what review or correction
is possible. A cryptographically valid action may be important evidence
of what a technical system accepted, but it does not by itself settle
the institutional legitimacy of the procedure or the authority
exercised.

\subsection{7.2 Shared capabilities and constituted
authority}\label{shared-capabilities-and-constituted-authority}

Organizations may reuse a model, software component, analytical method,
source collection, or agent configuration developed elsewhere. Such
reuse may reduce duplication and allow organizations to draw on
capabilities they did not build independently. Suitability, however,
remains context-dependent: a capability that is useful in one setting
may omit information, assumptions, or procedures needed in another.

The governing distinction is:

\begin{quote}
\textbf{Capabilities may be shared; institutional authority is not
transferred by technical reuse alone.}
\end{quote}

A reused capability does not automatically acquire the mandates of its
developer or the receiving organization. At the same time, the receiving
organization and any applicable common authority may legitimately
establish powers of representation, supervision, or coordination through
appropriate institutional instruments. Whether such powers exist, their
scope, and how they may be exercised depend on the relevant rules and
authority.

The distinction also applies to cognitive influence. A shared service
might possess no formal authority over a local decision while still
influencing it through source selection, default wording, availability
of alternatives, or control over updates. A technology-neutral
governance analysis should therefore examine both \textbf{formal
authority} and \textbf{practical control over institutional
intelligence}.

\subsection{7.3 Local discretion and shared
commitments}\label{local-discretion-and-shared-commitments}

Where a federation has applicable common commitments, local action may
need to be considered in relation to both local rules and shared
obligations. The resulting allocation of competence is an institutional
question; it cannot be deduced from technical architecture alone.

An authorized technical mechanism may implement a specific rule within
its established scope. For example, an institution might configure a
system to refuse an action when an expressly defined condition is not
met. Whether that refusal is authorized, contestable, or subject to
exception depends on the institutional arrangement.

An ambiguous question presents a different problem. An AI system may
identify an apparent conflict between a local proposal and a shared
commitment, but the interpretation and disposition of that conflict must
be attributed to the authority competent under the applicable rules. The
model does not presume that every such matter requires the same body,
procedure, or intervention.

Federation is consequently one setting in which authority, institutional
intelligence, and delegation interact. It is not a required form of
AI-assisted governance.

Appendix A illustrates this interaction in a fictional
multi-organization setting where autonomous institutions reuse
analytical capabilities and observe a constituted common reserve
commitment while retaining distinct local mandates and political
standing. Rather than collapsing local and federative decisions into a
single event, the appendix traces four linked Institutional Decision
Episodes---local adoption, federative detection and provisional
restriction, competent adjudication, and local reconsideration---to show
how shared cognitive tools, jurisdictional boundaries, and
contestability conditions can be analyzed without treating federation as
a universal requirement of the framework.

\section{8. Counterarguments, Risks, and
Limitations}\label{counterarguments-risks-and-limitations}

The framework's normative premise and conceptual distinctions do not
ensure that an institution will preserve meaningful human governance.
The following counterarguments identify conditions under which the
proposed arrangement may fail or produce effects contrary to its
objectives.

\subsection{8.1 Cognitive dependence and ceremonial
authority}\label{cognitive-dependence-and-ceremonial-authority}

AI-assisted analysis may make information easier to access, but it may
also lead participants to defer to generated outputs. Dependence could
arise where reviewers lack time, expertise, access to underlying
evidence, or practical authority to reject a recommendation.

A human approval step would then record a formal disposition while
leaving much of the practical framing and evaluation outside effective
human scrutiny. This is a challenge to the framework's human-authority
premise: authority retained on paper may not correspond to control
exercised in practice.

Human--AI interaction guidance and critiques of oversight support
examining these conditions; they do not establish that cognitive decline
or ceremonial approval inevitably follows every use of assistance
(\axCite{amershi2019guidelines}{Amershi et al.}; \axCite{green2022flaws}{Green}). The effects of assistance on
judgment and skill require context-specific evidence.

\subsection{8.2 Agenda capture and political
asymmetry}\label{agenda-capture-and-political-asymmetry}

AI may influence governance before a formal decision is made. A system
that selects records, recommends a problem definition, or offers a
limited set of options can affect which matters receive attention and
how alternatives are understood.

This influence may be concentrated in model providers, administrators,
source curators, institutional officers, or participants with better
access to private tools. Formal equality of voting rights may therefore
coexist with unequal capacity to shape deliberation.

The proposition is not that all AI assistance produces capture. It is
that a framework concerned only with formal decision rights could
overlook the institutional power exercised through agenda formation,
information selection, and control over cognitive resources.

\subsection{8.3 Epistemic monoculture and correlated
errors}\label{epistemic-monoculture-and-correlated-errors}

Several AI systems may appear to provide independent assessments while
relying on related models, shared source collections, similar
instructions, or common institutional assumptions. Under those
conditions, their errors may be correlated.

Shared foundations do not imply identical outputs or inevitable failure.
Conversely, technical diversity alone does not establish substantive
independence: systems built differently may still reproduce the same
omission or institutional bias.

An institution considering multiple forms of AI assistance should
therefore ask what kinds of dependence matter for the decision at hand
and whether participants can access alternative interpretations or
relevant primary evidence.

\subsection{8.4 Limits of adversarial
agents}\label{limits-of-adversarial-agents}

An adversarially assigned system may identify objections that a
proponent overlooked. Its label, however, does not prove impartiality,
independence, competence, or sufficient authority to make an objection
consequential.

An agent developed by the same team or using the same resources as a
compositional system may still produce useful criticism. It may also
share incentives or assumptions that narrow its challenge. The relevant
question is what the adversarial function actually examines, who
controls it, and whether its findings can be questioned or acted upon.

If an institution treats the existence of an adversarial agent as proof
that adequate scrutiny occurred, the procedure could provide assurance
without establishing that the underlying challenge was meaningful.

\subsection{8.5 False memory and interpretive
authority}\label{false-memory-and-interpretive-authority}

AI-assisted memory systems may use document retrieval, structured
databases, language models, or combinations of these methods. They can
consult records, but their usefulness depends on the selection,
currency, access conditions, and interpretation of those records.

A system may retrieve an obsolete policy, omit an amendment, misread the
scope of a precedent, or present an inferred relationship as though it
were an authoritative institutional conclusion. Participants may then
mistake a generated summary for the underlying record or for a binding
interpretation.

The framework therefore distinguishes retrieval of a record from a
model-generated interpretation of that record. Provenance can support
inspection, but it does not itself establish that the record is
complete, current, or legally authoritative (\axCite{w3c2013provdm}{W3C}).

\subsection{8.6 Institutional compliance and political
merit}\label{institutional-compliance-and-political-merit}

A proposal may satisfy the procedural rules applicable to it while
remaining undesirable, unfair, wasteful, or politically contested.
Equally, a proposal's apparent merit does not cure a defect in the
authority or procedure through which it is adopted.

AI-assisted analysis may help identify compliance questions and explain
possible consequences. It cannot turn procedural conformity into a
conclusive judgment about political merit. The distinction is especially
important when system-generated summaries present one course of action
as the natural or preferred result of a rule check.

\subsection{8.7 Procedural burden and informal
bypass}\label{procedural-burden-and-informal-bypass}

Evidence, review, and contestation can impose costs. If a framework
requires the same depth of documentation and review for every action,
participants may experience delay or workload disproportionate to the
consequences of the decision.

One possible institutional response is to bypass formal procedures;
another is that routine work simply becomes slower. These are risks and
testable possibilities, not established outcomes of the proposed model.

Proportionality is consequently part of the framework's design
objective. The appropriate level of evidence and review requires
institutional judgment and may vary with consequence, reversibility,
affected parties, and available capacity. Existing work on reviewability
informs the need to make processes reviewable; it does not by itself
demonstrate that a particular review regime causes paralysis (\axCite{cobbe2021reviewable}{Cobbe et al.}).

\subsection{8.8 Evidence access and meaningful
review}\label{evidence-access-and-meaningful-review}

Providing a large volume of records does not necessarily enable
effective challenge. Evidence may be late, incomplete, inaccessible,
technically difficult to interpret, or unavailable to the people
affected by a decision.

Reviewability should therefore be distinguished from practical
contestability. An institution may have records that permit a technical
audit while an affected participant lacks standing, assistance, time, or
an independent reviewer. Contestability research supports examining
intervention and challenge as institutional design concerns; the
adequacy of a particular procedure remains to be evaluated (\axCite{alfrink2023contestable}{Alfrink et al.}).

\subsection{8.9 Irreversibility and corrective
limits}\label{irreversibility-and-corrective-limits}

Some delegated actions can be suspended or corrected before their
effects become durable. Others may create consequences that cannot be
fully undone, even if an institution later identifies an authorization
or execution failure.

The framework distinguishes \textbf{revocation} of future authority,
\textbf{suspension} of activity, \textbf{correction} of a process or
record, \textbf{reversal} of an action where possible, and
\textbf{remedy} for consequences that remain. These concepts do not
establish that every remedy is available or adequate.

Where action is difficult to reverse, an institution may choose to
examine preventive review or narrower delegation. The appropriate choice
depends on the decision, the relevant rights and risks, and the rules
under which the institution operates.

\subsection{8.10 Jurisdictional and legitimacy
limits}\label{jurisdictional-and-legitimacy-limits}

The framework is a normative-conceptual account, not a universal legal
theory. Institutional authority, legal personhood, agency, standing,
contract effects, and liability may vary across jurisdictions.

An internal finding that an action exceeded a mandate does not
automatically determine the action's legal effect on third parties.
Likewise, a procedure may be institutionally recognized while the
legitimacy of the institution or procedure remains contested. Specific
legal consequences require the applicable legal context and competent
review.

\section{9. Research Propositions and Future
Evaluation}\label{research-propositions-and-future-evaluation}

The propositions below identify possible directions for future research,
not reported findings or authorization to conduct studies. Before
evaluation, specify the comparator, primary outcome, smallest relevant
effect, and results that would count against the hypothesis. Null,
adverse, and inconclusive findings must be distinguished. Any later
protocol requires appropriate access, institutional approval, and
relevant privacy, legal, and research-ethics review.

For each proposition, a null or adverse result would be informative.
Comparator conditions should receive equivalent evidence and time unless
access is itself the intervention; independent rubrics should also
capture false positives, omissions, uncertainty, and cases that do not
fit the scheme.

\subsection{9.1 Proposal composition}\label{proposal-composition}

\textbf{Hypothesis.} Compared with a standard template, source-linked AI
assistance will improve procedural completeness without increasing
material errors beyond a prespecified margin.

\textbf{Candidate observable outcomes.} Correct identification of
required information, omission of relevant alternatives, participant
drafting time, and errors introduced by the assistance.

\textbf{Possible comparison.} Participants using a standard template
could be compared with participants using AI assistance on equivalent,
synthetic proposal tasks.

\textbf{Confounders and limitations.} Prior expertise, language,
accessibility, tool familiarity, task design, and the quality of the
underlying rules may affect results. Procedural completeness does not
establish proposal merit.

\textbf{Evidence required.} A defined protocol, independent assessment
criteria, records of omissions and introduced errors, and reporting of
positive, null, and negative effects.

\textbf{Potential contrary result.} A null or adverse result, such as
equal completeness with materially more errors, would count against the
hypothesis.

\subsection{9.2 Cognitive access, understanding, and political
influence}\label{cognitive-access-understanding-and-political-influence}

\textbf{Hypothesis.} Source-linked assistance will improve comprehension
without reducing detection of erroneous recommendations. Influence over
agenda selection remains a separate research question.

\textbf{Candidate observable outcomes.} Comprehension of source
material, recognition of omitted alternatives, confidence calibration,
and differences in participants' opportunity to influence an agenda.

\textbf{Possible comparison.} Compare source-only, standard-summary, and
AI-assisted conditions using the same decision materials.

\textbf{Confounders and limitations.} Reading ability, prior knowledge,
trust, accessibility, source quality, and differences between simulated
influence and actual institutional power. Improved comprehension would
not by itself establish political equality.

\textbf{Evidence required.} Transparent measures, independent review of
framing, analysis across participant groups where appropriate, and
documentation of null or adverse effects.

\textbf{Potential contrary result.} A rise in confidence without
corresponding comprehension, or worse detection of incorrect advice,
would count against the hypothesis.

\subsection{9.3 Adversarial
contestation}\label{adversarial-contestation}

\textbf{Hypothesis.} Assisted critique will identify a greater
proportion of independently defined material issues without a material
increase in false objections.

\textbf{Candidate observable outcomes.} Relevant issues identified,
false objections, missed objections, participant workload, and
participants' willingness to challenge the adversarial analysis itself.

\textbf{Possible comparison.} Compare deliberation without an
adversarial brief, with a human-produced brief, and with an AI-assisted
brief, where such comparisons are feasible and ethically approved.

\textbf{Confounders and limitations.} Expertise, topic familiarity,
brief length, shared source material, model dependence, and
institutional incentives. A greater number of objections is not
necessarily better scrutiny.

\textbf{Evidence required.} Independently defined issue criteria,
assessment of false positives and omissions, and examination of whether
challenges lead to meaningful institutional consideration.

\textbf{Potential contrary result.} An increase in objection volume
without better identification of material issues, or with materially
worse precision, would count against the hypothesis.

\subsection{9.4 Institutional memory}\label{institutional-memory}

\textbf{Hypothesis.} Provenance-linked retrieval will improve
identification of operative records without increasing access or
disclosure errors.

\textbf{Candidate observable outcomes.} Retrieval precision,
identification of superseded records, accuracy of source interpretation,
time to locate relevant material, and disclosure or access errors.

\textbf{Possible comparison.} Compare conventional search and
AI-assisted retrieval on a controlled collection containing current,
amended, superseded, and irrelevant records.

\textbf{Confounders and limitations.} Quality and completeness of
archives, access permissions, changes in rules over time, and the
distinction between finding a record and determining its authority.

\textbf{Evidence required.} A documented source set, independently
verified status of records, error analysis, and privacy and access
controls appropriate to the material.

\textbf{Potential contrary result.} More obsolete-record selection or
access errors despite faster retrieval would count against the
hypothesis.

\subsection{9.5 Decision-episode reconstruction and failure
detection}\label{decision-episode-reconstruction-and-failure-detection}

\textbf{Hypothesis.} Relative to an evidence-matched case form,
IDE-based reconstruction will reduce unsupported inferences about
episode boundaries, MVA, MC, and execution.

The representation to be evaluated would use the accepted Stage 1
notation:

\[
\mathcal{E}=\langle P,R,M,C,T,A,E,Q,D\rangle,
\qquad MC(A,M,R).
\]

\textbf{Candidate observable outcomes.} Reconstruction accuracy, correct
identification of missing evidence, classification of the three
concerns, reviewer disagreement, time, and cases that do not fit the
categories.

\textbf{Possible comparison.} Independent reviewers could examine
equivalent synthetic records presented as conventional case files and as
episode-structured materials.

\textbf{Confounders and limitations.} Reviewer training, familiarity
with the notation, case complexity, quality of the underlying evidence,
and the risk that the structured format directs reviewers toward
predetermined categories.

\textbf{Evidence required.} Independently prepared cases and scoring
criteria, reporting of disagreements and unclassifiable cases, and
comparison of workload as well as accuracy. A faster or more consistent
classification would not by itself demonstrate institutional legitimacy
or effective remedy. In a separately authorized assessment, two
independent analysts should code the same fixed records without
conferring, compare episode boundaries and all nine fields, and
separately compare MVA and MC. A reviewer should document the basis of
any adjudication and preserve disagreements that the records cannot
resolve. No such independent test was conducted here.

\textbf{Potential contrary result.} An independently repeated result of
no reduction, or an increase in unsupported inferences, would count
against the hypothesis under the specified comparison.

\subsection{9.6 Delegation and practical
intervention}\label{delegation-and-practical-intervention}

\textbf{Hypothesis.} Explicitly distinguishing rules, mandate,
capability, enablement, and contestability will reduce false inferences
about authorization and practical intervention.

\textbf{Candidate observable outcomes.} Detection of mandate-scope
discrepancies, recognition of unavailable intervention, correct
distinction among suspension, revocation, reversal, and remedy, and
uncertainty reported by reviewers.

\textbf{Possible comparison.} Use hypothetical episodes in which mandate
validity, technical permission, action conformity, and review capacity
vary independently and in combination.

\textbf{Confounders and limitations.} Ambiguous rules, incomplete
evidence, institutional variation, and the artificial simplicity of
controlled scenarios. A hypothetical mismatch is not evidence that any
actual system contains the same defect.

\textbf{Evidence required.} Clearly stated assumptions, cases with
overlapping or uncertain failures, independent review of the
classifications, and documentation of situations in which the framework
adds no useful distinction.

\textbf{Potential contrary result.} Persisting confusion between
technical enablement and authority or between approval and external
execution would count against the hypothesis.

\subsection{9.7 Cross-institutional
applicability}\label{cross-institutional-applicability}

\textbf{Research question.} Which distinctions remain informative across
selected institutional forms, and which require adaptation or omit
important local features? No universal portability hypothesis is
assumed.

\textbf{Candidate observable outcomes.} Whether analysts can identify
authority, process, technical mechanisms, cognitive functions, and
relevant review conditions without omitting important
institution-specific features.

\textbf{Possible comparison.} Apply the framework to a small,
purposively selected set of institutional cases and compare its account
with alternative conceptual descriptions.

\textbf{Confounders and limitations.} Unequal access to institutional
records, jurisdictional differences, case selection, translation of
local terminology, and the possibility that the framework forces unlike
arrangements into misleading categories.

\textbf{Evidence required.} Explicit case-selection criteria, access and
confidentiality arrangements, attention to counterexamples and residual
categories, and a clear boundary between conceptual comparison and
generalizable empirical claims.

\textbf{Potential contrary result.} Recurring loss of
institution-specific features or adaptations that prevent comparison
would limit the applicability claim.

\section{10. Conclusion}\label{conclusion}

This paper proposes a bounded integrative conceptual framework for
examining how human institutions may incorporate AI into governance
while preserving legitimately constituted authority, meaningful
participation, accountability, and institutional autonomy. Constitution
alone does not establish normative legitimacy.

The framework distinguishes four substantive dimensions. \textbf{Human
legitimacy} concerns the institutional basis of decision rights and
authority. \textbf{Institutional intelligence} concerns the cognitive
work through which proposals are composed, evidence is considered,
claims are contested, and institutional memory is made available.
\textbf{Agentic delegation} concerns the relationship among mandates,
operational capability, technical permission, action, and
responsibility. \textbf{Federated institutional autonomy} concerns how
distinct institutions may organize their own capabilities while
observing applicable shared commitments; federation remains an
application context rather than a universal requirement.

The three analytical levels distinguish institutional authority,
governance processes, and technical mechanisms. The Institutional
Decision Episode connects these levels around a bounded action or
decision. Its separation of mandate validity, delegated-action
conformity, and contestability capacity is an integrative device for
examining different institutional concerns. The hypothetical procurement
episode illustrates how those concerns may lead to different questions
about evidence, responsibility, and possible responses; it does not
demonstrate that the device improves governance or determines legal
outcomes.

The framework also faces substantial counterarguments. AI assistance may
shape agendas, encourage dependence, reproduce shared errors, or present
generated interpretations as institutional authority. Formal human
approval, source records, or a named adversarial function do not
independently resolve those risks. Conversely, extensive review can
impose burdens that require proportionate institutional judgment.

The paper advances a conceptual and normative research direction with
potentially incremental diagnostic utility at institutional boundaries.
Its practical usefulness, empirical effects, inter-analyst
reproducibility, and applicability across institutional forms remain
open to evaluation. Its contribution is a synthesis of established
concerns rather than a claim of distinctive originality or demonstrated
superiority.

\begin{quote}
\textbf{Intelligence may be distributed; authority must be constituted.}
\end{quote}

\section{Author Contributions and
Responsibility}\label{author-contributions-and-responsibility}

The human author assumes intellectual responsibility for the research
presented in this manuscript, including its conceptual framing,
normative premises, interpretation of the literature, analytical
contributions, scientific limitations, and conclusions.

The manuscript was developed through a human-directed, AI-assisted
research process. Artificial intelligence systems supported research
organization, conceptual analysis, literature examination, manuscript
drafting, critical evaluation, bibliographic reconciliation, and
editorial preparation.

AI-assisted contributions were subject to human-directed scientific
review and decision-making processes. Such assistance does not confer
academic authorship or independent institutional authority upon the AI
systems involved. The available process record does not establish
independent verification of every AI-assisted output or cited source.

The human author remains responsible for the accuracy and integrity of
the submitted manuscript, appropriate scholarly attribution, critical
assessment of AI-assisted material, and final approval of any external
submission.

\section{Declaration of AI-Assisted Research and Manuscript
Development}\label{declaration-of-ai-assisted-research-and-manuscript-development}

Artificial intelligence systems were substantively involved in the
development of this conceptual research manuscript.

Their contributions included assistance with research structuring,
literature mapping, conceptual comparison, analytical argumentation,
manuscript drafting, critical examination, bibliographic reconciliation,
and editorial consolidation.

The research process was organized through human-directed
decision-making and scientific review, with AI systems functioning as
analytical and documentary collaborators rather than independent
academic authors.

The inclusion of AI-assisted analysis and text does not establish the
empirical validity, theoretical originality, or practical effectiveness
of the conceptual framework presented in this manuscript.

The human author retains responsibility for the manuscript's
intellectual content, scholarly integrity, interpretation, attribution,
and final publication decisions.

AI tools and models: OpenAI Codex assisted with pre-review bibliographic
and editorial reconciliation and preparation of this declaration. The
model version and the identities of AI systems used in earlier research
and drafting phases are not independently verified and remain to be
confirmed before external submission.

\phantomsection
\pdfbookmark[1]{References}{references-bookmark}
\label{references}
\bibliographystyle{axodus-authorial}
\bibliography{references}

\section{Appendix A. Illustrative Application: Federated Institutional
Governance and Contested Common
Commitments}\label{appendix-a.-illustrative-application-federated-institutional-governance-and-contested-common-commitments}

This appendix presents a fictional, nonempirical scenario illustrating
how independent human institutions can share analytical capabilities and
accept legitimately constituted common commitments without automatically
transferring local identity, political standing, mandate, or
institutional authority. It is an application of the paper's conceptual
framework and canonical Institutional Decision Episode representation,
not a new theoretical model, a universal constitutional prescription, or
a description of any operating institutional federation.

\subsection{A.1 Fictional federative setting and enumerated
competences}\label{a.1-fictional-federative-setting-and-enumerated-competences}

Consider three fictional, human-governed organizations---Alpha, Beta,
and Gamma---that have voluntarily acceded to a common federative compact
through procedures valid under their respective constitutive rules.
Their internal governance arrangements differ: Alpha decides major
treasury allocations through direct member voting; Beta governs through
an elected representative council; and Gamma combines member voting with
topic-specific delegation to human delegates. Neither the compact nor
this illustration assumes autonomous AI voting rights.

Under the hypothetical compact, each organization retains local autonomy
over project selection, internal governance, and treasury management,
subject to common obligations validly accepted through accession. The
federation possesses only expressly enumerated competences:
administering a shared settlement and liquidity facility and monitoring
compliance with a common reserve commitment designated \textbf{F-17}.
The federative administration may issue incompatibility notices and,
when defined conditions are met, impose a time-limited provisional
restriction limited strictly to use of the shared facility. A separately
constituted federative adjudication panel is empowered to interpret F-17
and determine whether a contested transaction is eligible to use the
shared facility. Neither federative body has authority to freeze a
member's local treasury generally, alter local membership standing,
annul a local vote, or substitute a replacement local budget allocation.

\subsection{A.2 Local institutional intelligence and capability
reuse}\label{a.2-local-institutional-intelligence-and-capability-reuse}

Within Organization Alpha, AI-assisted functions support proposal
composition, retrieval of local rules and prior resolutions, comparative
impact analysis, adversarial critique, and assembly of evidence for
challenges. These functions expand the analytical resources available to
participants, but they remain preparatory. They do not replace Alpha's
voting procedure, and generated summaries must remain distinguishable
from underlying institutional records.

Furthermore, formal local autonomy does not guarantee cognitive
independence. If Alpha, Beta, and Gamma rely on shared foundation models
or common retrieval indices to draft and critique proposals, their
deliberations may converge around similar framings, blind spots, or
correlated errors despite distinct voting rules. Assigning separate
labels such as ``composer'' and ``adversarial auditor'' does not by
itself establish independent judgment.

Suppose Alpha upgrades its internal critique workflow by adopting an
analytical module originally developed by Beta. Alpha configures the
module inside its own environment and defines its local advisory
mandate. Consistent with Section 7.2, \textbf{technical reuse of a
capability does not transfer institutional authority}. The imported
module carries neither Beta's institutional identity, credentials, and
political standing nor any authority for Beta to govern Alpha.
Conversely, if the module retains embedded policy defaults from
Beta---such as a rule requiring executive council sign-off---and Alpha's
configuration permits that rule to block a community-approved proposal,
Alpha experiences a delegated-action conformity failure caused by
unexamined policy leakage. Technical tenant boundaries and descriptive
capability or trait records do not by themselves establish political
autonomy, executable import permission, or freedom from inherited
assumptions.

\subsection{A.3 The ALPHA-42 / F-17 dispute and provisional restriction
requirements}\label{a.3-the-alpha-42-f-17-dispute-and-provisional-restriction-requirements}

Suppose Alpha's members deliberate on \textbf{ALPHA-42}, a fictional
proposal to deploy \textbf{230 reserve units} from Alpha's treasury into
an external yield facility routed through the shared federative
facility. An AI-assisted composition tool helps structure the proposal
text, while an adversarial review function flags a potential conflict
with \textbf{F-17}, which caps covered external encumbrances at
\textbf{30\% of eligible reserves}. After reviewing the proposal, the
objection, and a rebuttal arguing that the external position remains
liquid on demand, Alpha's members adopt ALPHA-42 under local rules. That
vote constitutes a local decision; it neither settles the interpretation
of F-17 nor executes the transfer.

On the stipulated valuation snapshot for the episode, Alpha holds
\textbf{1,000 eligible reserve units} (the fixed denominator for this
assessment) and already has \textbf{150 units} in covered external
encumbrances. If ALPHA-42's \textbf{230 units} are classified as a
covered encumbrance with no offsets, Alpha's covered total would rise to
\(150 + 230 = 380\) units, or \textbf{38\%} of eligible
reserves---exceeding the \textbf{300-unit (30\%)} ceiling by \textbf{80
units}. The arithmetic is determinate on the stipulated snapshot, but
whether the withdrawal terms make the position a covered ``encumbrance''
under F-17 is disputed.

A federative monitoring function detects the queued facility request and
reports a potential F-17 incompatibility. That technical flag is an
analytical observation, not a self-executing adjudication. Before any
provisional hold takes effect on the shared facility route, the compact
requires: (1) a competent decision-maker; (2) a defined trigger; (3)
scope restricted to the shared facility; (4) a maximum duration; (5)
prompt notice and disclosure of supporting evidence and valuation
records; (6) standing and an affordable, timely channel for Alpha to
challenge the hold; (7) review by a competent, sufficiently impartial
body free from disqualifying conflicts of interest; and (8) explicit
termination and remedy rules.

The federative administration verifies the trigger and places a
temporary hold on the shared facility route. Alpha challenges both the
hold and the classification of ALPHA-42 before the adjudication panel.
After examining the compact, the contract terms, and both parties'
submissions, the panel rules that the external facility's utilization
constraints create withdrawal risk covered by F-17. It therefore holds
the original 230-unit request ineligible to use the shared facility.
Because the panel lacks competence to rewrite Alpha's local budget, it
does not substitute a modified amount.

The matter returns to Alpha. Operating under its own local rules---and
subject to local notice, challenge, and reconsideration
procedures---Alpha's members vote on a revised proposal authorizing a
\textbf{130-unit} placement. On the same 1,000-unit snapshot, covered
encumbrances total \(150 + 130 = 280\) units (\textbf{28\%}). Following
Alpha's separate approval of the revised 130-unit proposal, any
subsequent transaction submission would require a distinct and valid
execution mandate, an authorized local executor, updated exposure and
valuation checks, and applicable clearance to use the shared facility.
Submission would constitute a separate operational action requiring its
own evidence and disposition. The local approval alone does not
establish that the transaction was submitted, accepted, or executed.

\subsection{A.4 Four linked Institutional Decision
Episodes}\label{a.4-four-linked-institutional-decision-episodes}

Rather than collapsing local and federative authority into a single
principal, the sequence can be examined as four linked episodes using
\(\mathcal{E}=\langle P,R,M,C,T,A,E,Q,D\rangle\) and \(MC(A,M,R)\):

{\def\LTcaptype{none} % do not increment counter
\setlength{\parskip}{0pt}\begin{longtable}[]{@{}
  >{\raggedright\arraybackslash}p{(\linewidth - 6\tabcolsep) * \real{0.1650}}
  >{\raggedright\arraybackslash}p{(\linewidth - 6\tabcolsep) * \real{0.2150}}
  >{\raggedright\arraybackslash}p{(\linewidth - 6\tabcolsep) * \real{0.3200}}
  >{\raggedright\arraybackslash}p{(\linewidth - 6\tabcolsep) * \real{0.3000}}@{}}
\toprule\noalign{}
Episode & Principal, Rules, and Mandate (\(P,R,M\)) & Capability, Permission, and Action (\(C,T,A\)) & Evidence, Contestability, and Disposition (\(E,Q,D\)) \\
\midrule\noalign{}
\endhead
\bottomrule\noalign{}
\endlastfoot
\textbf{IDE-B1 --- Local Decision} & \(P\): Alpha voting members. \(R\):
Alpha bylaws and accepted compact obligations. \(M\): local competence
to approve treasury proposals. & \(C\): AI drafting/critique and voting
system. \(T\): a distinct voting or recording enablement mechanism is
not evidenced in the examined narrative. \(A\): adoption of ALPHA-42
(230 units). & \(E\): proposal drafts, audit warning, vote tally. \(Q\):
local procedural challenge. \(D\): conditional local approval; \(MC\) is
locally conforming, while F-17 coverage remains contested. \\
\textbf{IDE-B2 --- Federative Detection \& Provisional Action} & \(P\):
federative administration. \(R\): compact notice and provisional-hold
rules. \(M\): bounded authority over shared-facility access. & \(C\):
cross-facility screening. \(T\): the stipulated route hold was applied;
a distinct configured enablement mechanism is not evidenced in the
examined narrative. \(A\): incompatibility notice and temporary facility
hold. & \(E\): F-17 revision, 1,000-unit snapshot, calculation
(380/1,000), notice timestamp. \(Q\): expedited challenge and expiry
review. \(D\): temporary route hold; \(MC\) depends on trigger, scope,
and duration compliance. \\
\textbf{IDE-B3 --- Federative Adjudication} & \(P\): constituted
adjudication panel. \(R\): compact dispute-resolution rules. \(M\):
authority to interpret F-17 and rule on facility eligibility. & \(C\):
documentary analysis, evidentiary evaluation and reasoned adjudication.
\(T\): a distinct configured ruling-registration mechanism is not
evidenced; adjudicative competence does not itself establish such
enablement. The panel has no stipulated competence to operate the shared
facility or modify Alpha's allocation; administrative actors implement
restrictions or clearances. \(A\): adjudication of F-17 coverage and
determination of the original 230-unit request's eligibility to use the
shared federative facility. & \(E\): party briefs, contract terms,
hearing record. \(Q\): challenge to panel competence or impartiality
under compact rules. \(D\): 230-unit request ruled ineligible for the
shared facility; no local budget substitution. \\
\textbf{IDE-B4 --- Local Response} & \(P\): Alpha's competent local
voting body. \(R\): local reconsideration rules and applicable common
obligations. \(M\): local competence to approve the revised allocation.
& \(C\): proposal preparation, deliberation and voting capabilities.
\(T\): a distinct voting or recording enablement mechanism is not
evidenced in the examined narrative. External-transfer enablement is not
applicable to this local approval; any such enablement for a later
operation is not evidenced here. \(A\): adoption of the revised 130-unit
proposal. & \(E\): revised proposal, deliberation record, stipulated
valuation assumptions and local voting record. \(Q\): local notice,
challenge and reconsideration procedures. \(D\): revised local approval;
no facility clearance, submission, acceptance, or external execution is
established by the examined B4 record. These would belong to a
subsequent operational action. \\
\end{longtable}
}

For B4, the local body's competence to reconsider is stipulated by the
fictional compact and local arrangements. MVA may be treated as valid
for that competence under those assumptions; the full procedural
regularity of the vote remains undetermined from the stated record. The
revised local approval is compatible with the described deliberative
scope under the stipulated interpretation, while overall institutional
conformity remains conditional on evidence not supplied. The calculation
of 280/1,000 = 28\% demonstrates only arithmetic compatibility with the
stipulated 30\% cap at that snapshot; it does not establish mandate
validity, procedural regularity, future clearance, or execution. No
distinct procedural enablement mechanism follows solely from the vote.

Across these episodes, all three failure classes remain possible:
federative overreach beyond enumerated powers
(AUTHORITY\_VALIDITY\_FAILURE), unexamined policy leakage from a reused
module (DELEGATED\_ACTION\_CONFORMITY\_FAILURE), or prohibitive filing
costs, delays, and biased review that render appeal nominal
(CONTESTABILITY\_CAPACITY\_FAILURE). The illustration shows how shared
intelligence and common commitments can be analyzed across linked
jurisdictions without implying empirical effectiveness, frictionless
coordination, or an implemented cross-domain software stack.

\end{document}
